% =============================================================================================
% §8 Reproducibility and verification.
%
% The section's job is to let a reader decide how much to trust the numbers, which means being
% specific about which guarantees hold and which do not. A blanket "our results are reproducible"
% is worth less than a scoped claim a reader can check, so the section leads with the boundary
% rather than with the guarantee.
% =============================================================================================
\section{Reproducibility and verification}
\label{sec:repro}

\subsection{What is bit-exact, and what is not}
\label{sec:repro-bitexact}

Reproducibility here is a deliverable rather than hygiene, and it is scoped rather than asserted
whole. Three things are bit-exact from one configuration file, a pinned seed and content-hashed
inputs: the data pipeline, the frozen normalization statistics, and prediction replay. Any recorded
prediction can be reconstructed from the recorded inputs.

GPU training is not bit-exact by default, and the honest version of that sentence took us a
correction to write. The production path now sets deterministic algorithms explicitly, at a
measured throughput cost of $13.175\%$, and every run records the mode it ran in.
% artifact: docs/BUILD_RECORD.md §2 D20 — fixed at 591ec44, deterministic_algorithms=True on the production path; penalty 13.175% per docs/m6_cuda_probe_report.md

The correction is worth stating because it is the kind of error a reproducibility section normally
hides. For a period, $49$ of $49$ run records carried an assertion of bit-exactness beside a
configuration that had determinism switched off. The cause was a single line that computed the
claim as \emph{attention mode equals the deterministic kernel} --- our own invariant's sufficiency
premise encoded as code. Deterministic attention is necessary for bit-exact training and is not
sufficient, because reduction order elsewhere in the backward pass is unconstrained. The scope of
the false claim was established at the time and is narrow: only the GPU-training clause was
affected; the pipeline, frozen-statistics and prediction-replay clause held, the anchor of
\cref{sec:repro-anchor} held, and no prior result was invalidated. We report it here rather than
quietly shipping the fix because a reader assessing our other claims should know that this one was
wrong for a while and how it was caught.
% artifact: docs/BUILD_RECORD.md §2 D20 — "49 of 49 run records carried bit_exact_claim: true beside deterministic_algorithms=False"; root cause attention_mode.py `claim = attention_mode == MODE_SDPA`

One consequence survives the fix. Same-seed runs were observed to diverge --- fraction-negative
$0.5662$ against $0.99883$ for the same cell at identical horizon, cap, estimator, seed and recipe
--- which is basin-hopping, and a lower bound on across-seed variance. Forcing determinism removes
the same-seed component and leaves the across-seed component untouched, so it does not repair the
statistical consequence. \Cref{sec:lim-power} carries that; \cref{sec:repro-guards} describes what
we do about it.

\subsection{The anchor}
\label{sec:repro-anchor}

A reproducibility claim that is only checked at the end is a claim about one moment. We instead
carry a standing anchor: a fixed evaluation of a frozen checkpoint through the full harness, whose
results hash is re-proven before and after any change to the evaluation instrument. The current
anchor is \texttt{sha256:3f86882a63dd06c78\ldots}, over inputs that are themselves content-hashed
--- dataset, tokenizer and predictor each carry their own SHA-256 --- and it is accompanied by an
exhaustive causal sweep at $420$ of $420$ checks.
% artifact: runs_manifest/gate_a_run_manifest.json :: results_hash = sha256:3f86882a63dd06c780e7d73f61a5253e39f5b07564d4c84152bbaad62c886dc3; gate_a_causal_sweep.{coverage=420, total=420, passed=true}; inputs.{dataset_hash, tokenizer_hash, predictor_hash}; env.{python 3.11.7, numpy 2.4.6, torch 2.12.1}

Two procedural rules attach to it, both of which exist because of a specific failure. The anchor is
re-proven \emph{twice}, from separate invocations. And its exit status is read from an unpiped
command, because an exit code read through a pipe is the pipe's --- a `\texttt{| grep | tail}`
once masked a non-zero exit and let a false anchor claim into a commit message. The anchor has been
legitimately re-based once, when a file inside the anchored instrument changed for a correctness
fix; the re-basing is recorded rather than presented as continuity.
% artifact: docs/BUILD_RECORD.md §5 standing rules — pipefail rider (born from D13), anchor rule; the 5eead7b6 -> 3f86882a re-anchoring

The environment is pinned by a lockfile: Python 3.11.7, NumPy 2.4.6, PyTorch 2.12.1. That pin is
itself a response to a defect --- environment drift with no lockfile --- and load-bearing results
are re-verified under it rather than under whatever interpreter happened to be first on the path.
% artifact: uv.lock; pyproject.toml requires-python >= 3.11; docs/BUILD_RECORD.md §2 D5

\subsection{Content addressing}
\label{sec:repro-content}

Every produced dataset is content-hashed, and the universe carries a single Merkle root over the
lake, recorded in the manifest with per-symbol listing and delisting dates
(\cref{sec:data-provenance}). The property this buys is specific: any past prediction can be
reconstructed, because the exact bytes it was computed from are identified rather than described.
The compaction of the lake was proven content-preserving by re-deriving all $200$ per-symbol hashes
and recomputing the root, rather than by inspecting the compactor.

\subsection{Pre-registration, and an amendment log that records the direction}
\label{sec:repro-prereg}

The decision surface --- clauses, thresholds, horizon, cost, the $\kappa$ grid, the seed count and
the trial count --- was fixed before any of its inputs existed, and the design was frozen before
the run. Amendments are not forbidden; unlogged amendments are. Every change is recorded with its
date, its reason and, crucially, a \emph{direction-blind} test: before looking at whether the
amendment helps, we record whether it tightens or loosens the criterion. The log carries $41$
distinct revision tags.

The direction test is not decorative, and it has already overturned our own summary of our own
amendments. The most consequential window changed the multiple-testing trial count and the
dispersion basis for the deflation benchmark --- moved from all arms to the placebo arm alone
--- together, and it was written up as \emph{every
amendment in this window loosens the effective bar}. That was false. The deflation clause does
loosen, by a factor of $0.859$. But the clause that most directly tests the headline contrast
\emph{tightens}, twice over: its quantile multiplier rises from $2.4865$ to $3.0728$ at five seeds,
a factor of $1.236$, and the total standard error gains a training-variance term that can only
increase it. The net across a conjunctive gate is indeterminate before the run and is not
monotonically loosening. Had the original framing reached this paper it would have been a false
statement about our own amendments, in the section a referee reads most carefully. We report the
aggregate as one number rather than letting a favourable half stand for the whole.
% artifact: docs/m6_prereg.md §7 v1.5 — the amendment window, the direction-blind test per amendment, the aggregate-as-one-number safeguard

Where a specification change means two defensible analyses exist, both are computed from the same
artifacts and both are reported. \texttt{dual\_specification} is a required field of the verdict
manifest, and a manifest that lacks it cannot be quoted; disagreement between the two
specifications is a first-class finding rather than a footnote.
% artifact: src/trikaal/eval/verdict.py :: dual_specification assembled and emitted into the manifest; §7 v1.5 makes it a REQUIRED field

\subsection{The conformance surface}
\label{sec:repro-conformance}

Pre-registration is only binding if something checks it. Every pinned quantity is a named constant
that is hard-asserted before any cell is scored: the symbol set by hash, the $\kappa$ grid, the
five seeds, the flat headline cost of $0.0030$, the evaluation sequence length, the realized
backbone parameter count of $21{,}301{,}248$, the $\hat{\mu}$ estimator, the training step counts,
the bootstrap configuration, the causal-encoder flag, the pointwise-fine flag and the
microstructure class weight. A run whose configuration differs from any of these does not produce a
downgraded result; it raises and produces none.
% artifact: src/trikaal/eval/conformance.py :: PINNED_SYMBOLS_SHA256, PINNED_KAPPAS, PINNED_SEEDS, PINNED_HEADLINE_COST = 0.0030, PINNED_MONEY_SEQ_LEN = 512, PINNED_BACKBONE_PARAMS = 21_301_248, PINNED_MU_ESTIMATOR = "expectation", PINNED_STEPS_STAGE1/2 = 26_003, PINNED_BOOT = {B: 10000, seed: 20260704, alpha: 0.05}, PINNED_ENCODER_CAUSAL, PINNED_FINE_POINTWISE, PINNED_MICRO_POINT_WEIGHT = 3.0

\subsection{Guards that refuse rather than decide}
\label{sec:repro-guards}

Two guards sit outside the decision rule, and their design constraint is that neither can change an
outcome --- each can only withhold one. The degeneracy guard refuses a verdict for any cell whose
book is constant-direction. The power guard refuses one when a cell's across-seed range reaches the
size of the effect being claimed. Both are computed after the clauses, neither mutates them, and
neither can flip a survival to a null or the reverse.

That asymmetry is what made it legitimate to add them after the design was frozen: a guard that can
only produce \textsc{halt\_adjudicate} can manufacture neither a positive result nor a negative
one. It is also why they are described here rather than in \cref{sec:setup} --- they are part of
how the result is verified, not part of what is being tested. \Cref{sec:setup} states one
asymmetry we chose not to remove and the reason.
% artifact: src/trikaal/eval/verdict.py :: degeneracy_guard, power_guard — both HALT-only, computed post-hoc; tests/eval/test_degeneracy_guard.py includes a KAT proving byte-identical clauses and primary with and without a degenerate cell

\subsection{How the defects were actually found}
\label{sec:repro-defects}

We keep a defect ledger, and we report it because the distribution in it is more informative than
any assurance we could offer. It records $36$ defects: $19$ we classify as critical --- meaning
they would have produced a wrong published result --- $15$ high, and $2$ medium. All but one are
closed; the exception is the capacity-eviction finding itself, which is mitigated rather than
closed because it is a property of the objective and not a bug.

The useful column is not severity but \emph{how each was found}:

\begin{center}
\begin{tabular}{lrr}
\toprule
detection layer & defects & of which critical \\
\midrule
Independent audit (first pass) & 9 & 6 \\
Builder or supervisor, outside a formal pass & 8 & 3 \\
Independent re-audit & 6 & 2 \\
Execution (found by running it) & 5 & 2 \\
Design review & 4 & 2 \\
Planted-signal canary & 3 & 3 \\
Independent audit (third pass) & 1 & 1 \\
\midrule
Total & 36 & 19 \\
\bottomrule
\end{tabular}
\end{center}
% artifact: docs/BUILD_RECORD.md §2 — the ledger; layer and severity columns tabulated directly from it

\textbf{No single layer found even a third of the critical defects.} The largest contributor, the
first independent audit, accounts for $6$ of $19$. Three critical defects were found only by
planting a signal of known information content and observing that it was not recovered --- no
review of that code would have found them, because the code was correct and the \emph{interface}
was not. Two were found only by executing the pipeline, which is the basis for the operational rule
that a component whose first execution is on rented hardware is an untested component whatever its
coverage says.

The reason to publish this table is that it is the strongest available argument against the
verification story a paper normally tells. Code review, a green test suite and a careful author
would have caught roughly a third of what mattered here. We do not claim the remaining defects are
gone; we claim that seven independent detection layers were run, that each found something the
others missed, and that the count of defects found by the most recent layer has not yet reached
zero.

A recurring pattern is worth naming, because it changes what a reader should count as evidence.
Several defects were checks that could not fail: a mutation test whose fixture made the two cases
it separated numerically identical; a gate that shipped below a \texttt{return} and had never
executed; a unit test whose written contract was the defect. In each, a test existed and was
green. Our operating rule is therefore that a gate or negative test is suspect until it has been
observed to fail, and that closing a finding requires restoring the pre-fix source in an isolated
tree and watching the new test fail there.
% artifact: docs/BUILD_RECORD.md §5 standing rules — fixture-discrimination, "a check that cannot fail is not a check", remediation rule; scripts/m6_reaudit_mutations.py is the harness, with its own negative control

\subsection{What is released, and what a reader can check without it}
\label{sec:repro-release}

The artifact is the code, the trained weights, the pre-registration with its full amendment log,
the run manifests, and the scripts that regenerate every figure in this paper from those manifests
at render time. No number in any figure is transcribed; each is loaded from the receipt that
produced it, so a figure that disagrees with its artifact fails to build rather than rendering a
stale value.

Two things a reader can check without rerunning anything: the anchor hash, which fixes the
evaluation instrument, and \cref{app:claims}, which separates every claim in this paper into
measured and estimated, with the artifact path for each. \Cref{sec:limitations} states what remains
unestablished.
